\section{Выводы}
\begingroup
\raggedright

Чем больше ключей было в наборе, тем дольше работала программа и тем больше оперативной памяти ей требовалось на пике. Самой долгой оказалась смешанная нагрузка на 100\,000 ключах с общим префиксом. Её медиана составила \(535{,}552\)~мс. Для вставки диапазоны времени на случайных ключах и на ключах с общим префиксом пересекались при всех трёх размерах набора. На тысяче ключей медиана для общего префикса была ниже, на двух больших наборах выше. Устойчивого изменения в этих замерах не видно.

Даже в серии успешных поисков в большинстве выборок встречалось чтение команд. Чтобы измерить время самого поиска, его пришлось бы замерять отдельно от ввода и начального заполнения словаря.

На проверенных входах ответы совпали с ожидаемыми. AddressSanitizer не сообщил об ошибках доступа к памяти, а \texttt{leaks} не обнаружил утечек в своём сценарии. Исправлений по результатам этих проверок не потребовалось, поэтому сравнение версий не проводилось.
\par
\endgroup
